\section{On the naturalness of $\beta$}
\label{sec:beta_origin}

The DEV phenomenology hinges on a single dimensionless coupling
$\beta \equiv \gamma^{2}/(m_{A}^{2}\,a_{0})$, which our SPARC fit pins to
$\beta_{\rm best} = 0.0075$. A natural question is whether this number can be
derived from fundamental constants, or whether it must be regarded as an
independent input to the theory. We address this question with a deliberately
brutal numerical sweep, and report the conclusion honestly: there is
\emph{no} unforced dimensional combination that lands on $\beta_{\rm best}$
within better than $\sim 15\%$, and the closest candidate is at the level of
a numerical coincidence rather than a derivation.

\paragraph{(1) Dimensional sweep.} We tested fifteen combinations built from
$\rho_{\Lambda}$, $\rho_{\rm crit}$, $a_{0}$, $c$, $H_{0}$, $G$, and the DEV
vacuum density $\rho_{\rm vac}^{\rm DEV} = a_{0}^{2}/(16\pi G)$. The closest
match is $(\rho_{\Lambda}/\rho_{\rm crit})^{1/3} = 8.59\times 10^{-3}$, only
$15\%$ above $\beta_{\rm best}$ ($|\log_{10}{\rm ratio}| = 0.06$). Two other
candidates fall within an order of magnitude
($3\sqrt{\rho_{\Lambda}/\rho_{\rm crit}}$, $\sqrt{3\rho_{\Lambda}/\rho_{\rm crit}}$),
but no combination is exact, and the cube-root form has no obvious
field-theoretic motivation.

\paragraph{(2) Factor-of-nine puzzle.} Since
$\sqrt{\rho_{\Lambda}/\rho_{\rm crit}} \simeq 8.0\times 10^{-4}$, the empirical
$\beta_{\rm best}$ exceeds it by a factor $\simeq 9.4$. Searching for an
integer-DOF origin, $N_{\rm DOF}^{2}\sqrt{\rho_{\Lambda}/\rho_{\rm crit}}$ with
$N=3$ gives $7.16\times 10^{-3}$ (ratio $0.955$). This is suggestive, but
neither $4\pi/\sqrt{4\pi/3}$, $1/(2\alpha)$, nor $\Omega_{m}/\Omega_{\Lambda}^{2}$
reproduce $9$, and we resist over-interpreting an unexplained $N=3$.

\paragraph{(3) Scale consistency.} A weighted linear fit of
$\beta_{\rm implied}$ (extracted from $\eta(g)$ across five physical regimes
spanning $g/a_{0}\in[0.003,1]$) versus $\log_{10}(g/a_{0})$ yields slope
$b = -0.005 \pm 0.018$, consistent with zero at $0.26\sigma$. $\beta$ is
therefore scale-invariant within current uncertainties, as required for it
to be a constant of the theory.

\paragraph{(4) Coincidence pairs.} Sweeping $C_{1}^{\alpha}C_{2}^{\gamma}$ with
$C_{1}=a_{0}/(cH_{0})$ and $C_{2}=\sqrt{\rho_{\Lambda}/\rho_{\rm crit}}$ over
half-integer exponents in $[-2,2]$, three pairs land within a factor of two
of $\beta_{\rm best}$, the closest being $C_{1}\cdot C_{2}^{1/2}$ (ratio
$0.66$). None are exact.

\paragraph{(5) First-principles $\gamma$.} Adopting the lower bound
$m_{A} = 3.7\times 10^{-25}$ eV$/c^{2}$, the inversion
$\gamma = \sqrt{\beta\, m_{A}^{2}\, a_{0}}$ gives, when interpreted as a
dimensionless coupling in natural units, $\gamma_{\rm nat}\simeq\sqrt{\beta}=0.087$.

\paragraph{(6) Dark-photon comparison.} Identifying $\gamma\sim\varepsilon\, e$
with the dark-photon kinetic-mixing parameter requires
$\varepsilon\sim 9\times 10^{-2}$, which exceeds astrophysical and laboratory
bounds ($\varepsilon\lesssim 10^{-3}$) by nearly two orders of magnitude.
DEV $\gamma$ is therefore \emph{not} a generic kinetic mixing of the visible
photon with a dark $U(1)$ — it is a coupling internal to the DEV sector.

\paragraph{(7) Status.} In light of the above, we adopt the conservative
position that $\beta$ is, like the fine-structure constant $\alpha$, an
empirical input of the theory. The numerical proximity to
$(\rho_{\Lambda}/\rho_{\rm crit})^{1/3}$ and to $9\sqrt{\rho_{\Lambda}/\rho_{\rm crit}}$
is recorded as a coincidence worth revisiting once the DEV Lagrangian is
embedded in a UV completion, but we explicitly do \emph{not} claim a
derivation. The predictive power of the theory rests on the fact that this
single number, once fit to galactic data, controls the cluster, UDG, and
$f\sigma_{8}$ phenomenology without re-tuning.
